%!TEX root = ../../thesis-trier.tex


In~\cref{constrsync:sec:constraut:startfinal} we show that, in a constraint automaton, 
we have some degree of freedom on how to distribute the start state and the final states. 
Later, in some proofs, this allows us to choose those states 
in a way that fits best in a given situation.
%This comes in handy, as
%it allows us to choose a set of final states that fits best in a given situation.
In~\cref{constrsync:sec:constraut:concatlowerbounds} 
we remark how to reformulate some results from~\cref{constr_sync:subsec:concat} 
in terms of PDFAs.
Various conditions on constraint PDFAs implying
that the problem is equivalent to $\constrsync{\Sigma^*}$
are collected in~\cref{constrsync:sec:constraut:unconstrained}. 


\subsection{Distributing the Start State and the Final States}
\label{constrsync:sec:constraut:startfinal}


By~\cref{constr_sync:cor:factors_sets_included,gen_results:cor:factors_etc_complexity} the complexity
of the constrained synchronization problem depends on the set of factors of the constraint 
language. \Cref{constr_sync:prop:final_states}
concerns how we can distribute a start state and a set
of final states in a trim PDFA without changing the factor set of the accepted language.
Various conclusions from \cref{constr_sync:prop:final_states} 
on altering the start state and set of final states
while preserving the computational complexity are collected in \cref{constr_sync:move_states}.
Lastly, in \cref{constr_sync:prop:start_states}, 
we prove that for trim PDFAs, considering multiple starting
states does not change the complexity.

%Recall the definition of the reachability quasiorder $\lereach$ and the induced partial order on the strongly connected components
%from~\cref{prelim:sec:reach} which extends without difficulty to PDFAs.

\begin{proposition}
	\label{constr_sync:prop:final_states}
	Suppose $\automat{B}{P}{\Sigma}{\mu}{p_0}{F}$ is a trim PDFA. Let $p \in P$
	and $E \subseteq P$ be such that
	$p_0$ is accessible\footnote{As $\mathcal B$
		is assumed to be trim,  $p \lereach p_0$, i.e., $p_0$ being accessible in $\mathcal B_{p, E}$, is equivalent ot $p_0 \eqreach p$.}
	in $\mathcal B_{p, E}$
	and every state in $F$ is coaccessible in $\mathcal B_{p, E}$.
	Then $\fact{L(\mathcal B)} = \fact{L(\mathcal B_{p, E})}$.
\end{proposition}
\begin{proof} % maximal bedingung weg, nur corollar; und trim annahme weg?
	Let $u \in L(\mathcal B)$. 
	As $p_0$ is accessible in $\mathcal B_{p, E}$,
	there exists $v \in \Sigma^*$ such that 	
	$\applytrans{\mu}{p}{v} = p_0$.
	Furthermore, as $\applytrans{\mu}{p_0}{u} \in F$,
	and each state in $F$ is coaccessible in $\mathcal B_{p,E}$ 
	there exists $w \in \Sigma^*$
	such that $\applytrans{\mu}{\applytrans{\mu}{p_0}{u}}{w} \in E$.
	Hence $vuw \in L(\mathcal B_{p, E})$.
	
	Now, suppose $u \in L(\mathcal B_{p, E})$.
	As $\mathcal B$ is trim,
	there exist $v, w \in \Sigma^*$
	such that $\applytrans{\mu}{p_0}{v} = p$
	and $\applytrans{\mu}{\applytrans{\mu}{p}{u}}{w} \in F$.
	So $vuw \in L(\mathcal B)$.
\end{proof}

Next, using \cref{constr_sync:cor:factors_sets_included,gen_results:cor:factors_etc_complexity}, we state some consequences of~\cref{constr_sync:prop:final_states} regarding changes to the start state or modifications to the final state set, while preserving the computational complexity of the associated constrained problem.

\begin{corollary}
	\label{constr_sync:move_states}
	Suppose $\automat{B}{P}{\Sigma}{\mu}{p_0}{F}$
	is a trim PDFA.
	\begin{enumerate} 		
		\item \label{constr_sync:move_states:1} $\constrsync{L(\mathcal B)} \eqACZEROm 
		\constrsync{L(\mathcal B_{s, F \cup \{q\}})}$
		for every $s \in [p_0]_{\eqreach}$
		and $q \in P$.\footnote{Note that, as $\mathcal B$ is trim, we have $F \ne \emptyset$ and~$p$ is coaccessible.}
		
		
		\item \label{constr_sync:move_states:2} For $q \in P$ coaccessible in $\mathcal B_{p_0, F\setminus\{q\}}$ 
		we have $\constrsync{L(\mathcal B)} \eqACZEROm 
		\constrsync{L(\mathcal B_{p_0, F \setminus \{q\}})}$.	
		
		\item \label{constr_sync:move_states:3} $\constrsync{L(\mathcal B)} \eqACZEROm  \constrsync{L(\mathcal B_{p_0, P})}$.
		
		\item \label{constr_sync:move_states:4} Let $E \subseteq P$ be a set such that $E \cap [q]_{\eqreach} \ne \emptyset$ for each
		maximal strongly connected component $[q]_{\eqreach}$.
		Then $\constrsync{L(\mathcal B)} \eqACZEROm 
		\constrsync{L(\mathcal B_{p_0, E})}$.
		
	\end{enumerate} 
\end{corollary}
\begin{proof} In the following we only show equality of the factor sets,
	the equivalence for $\leACZEROm$-reductions follows from~\cref{gen_results:cor:factors_etc_complexity}.	
		
\ref{constr_sync:move_states:1}: By~\cref{constr_sync:prop:final_states} applied to $p = p_0$
	and $E = F \cup \{q\}$.	
	
\ref{constr_sync:move_states:2}:   By~\cref{constr_sync:prop:final_states} applied to $p = p_0$	and $E = F \setminus \{q\}$.
	
	
\ref{constr_sync:move_states:3}:  By~\cref{constr_sync:prop:final_states} applied to $p = p_0$
	and $P$ where, as $F \subseteq P$, every state in $F$ is coaccessible in $L(\mathcal B_{p_0, P})$.
	
\ref{constr_sync:move_states:4}: By the assumption
 every state is coaccessible in $B_{p_0, E}$
 so the claim follows from~\cref{constr_sync:prop:final_states}.
\end{proof}

	\begin{figure}[htbp]
	\centering	
	\begin{minipage}{0.2\textwidth}
		\begin{tikzpicture}[>=stealth,shorten >=1pt,auto, node distance = 1.8cm]
			\tikzset{every state/.style={minimum size=1pt,fill=white}}
			
			\node[state, initial above] (q0) {};
			\node[state, right = of q0] (q1) {};
			\node[state, below = of q1] (q2) {};
			\node[state, accepting, left  = of q2] (q3) {};
			
			\path[->] (q0) edge[bend left] node {$a$} (q1)
			(q1) edge[bend left] node {$b$} (q0)
			(q1) edge node {$a$} (q2);
			
			\path[->] (q2) edge [bend left] node {$a$} (q3)
			(q3) edge [bend left] node {$b$} (q2);
		\end{tikzpicture}
	\end{minipage}\hfill
	\begin{minipage}{0.2\textwidth}
		\begin{tikzpicture}[>=stealth,shorten >=1pt,auto, node distance = 1.8cm]
			\tikzset{every state/.style={minimum size=1pt,fill=white}}
			
			\node[state] (q0) {};
			\node[state, initial above, accepting, right = of q0] (q1) {};
			\node[state, below = of q1] (q2) {};
			\node[state, accepting, left  = of q2] (q3) {};
			
			\path[->] (q0) edge[bend left] node {$a$} (q1)
			(q1) edge[bend left] node {$b$} (q0)
			(q1) edge node {$a$} (q2);
			
			\path[->] (q2) edge [bend left] node {$a$} (q3)
			(q3) edge [bend left] node {$b$} (q2);
		\end{tikzpicture}
	\end{minipage}\hfill
	\begin{minipage}{0.2\textwidth}
		\begin{tikzpicture}[>=stealth,shorten >=1pt,auto, node distance = 1.8cm]
			\tikzset{every state/.style={minimum size=1pt,fill=white}}
			
			\node[state, initial above] (q0) {};
			\node[state, right = of q0] (q1) {};
			\node[state, accepting, below = of q1] (q2) {};
			\node[state, left  = of q2] (q3) {};
			
			\path[->] (q0) edge[bend left] node {$a$} (q1)
			(q1) edge[bend left] node {$b$} (q0)
			(q1) edge node {$a$} (q2);
			
			\path[->] (q2) edge [bend left] node {$a$} (q3)
			(q3) edge [bend left] node {$b$} (q2);
		\end{tikzpicture}
	\end{minipage}\hfill
	\begin{minipage}{0.2\textwidth}
		\begin{tikzpicture}[>=stealth,shorten >=1pt,auto, node distance = 1.8cm]
			\tikzset{every state/.style={minimum size=1pt,fill=white}}
			
			\node[state, accepting, initial above] (q0) {};
			\node[state, accepting, right = of q0] (q1) {};
			\node[state, accepting, below = of q1] (q2) {};
			\node[state, accepting, left  = of q2] (q3) {};
			
			\path[->] (q0) edge[bend left] node {$a$} (q1)
			(q1) edge[bend left] node {$b$} (q0)
			(q1) edge node {$a$} (q2);
			
			\path[->] (q2) edge [bend left] node {$a$} (q3)
			(q3) edge [bend left] node {$b$} (q2);
		\end{tikzpicture}
	\end{minipage}\hfill
	\caption[Distributions of final states that give equivalent constrained problems]{The following distributions
		of final states (marked with a double circle) and start state
		all yield the same complexity in the constrained synchronization problem
		when the accepted languages are used as constraint languages.}
	\label{constrsync:fig:distributions-final-states}
\end{figure}	

To illustrate \cref{constr_sync:move_states}, 
\cref{constrsync:fig:distributions-final-states} gives different distributions of final states and a start state
for a PDFA that all yield the same computational complexity in the constrained synchronization problem
when the constraint language is the language accepted by the PDFA.

When we omit the assumption of $\mathcal B$ being trim
in \cref{constr_sync:move_states}, \ref{constr_sync:move_states:1} still holds
for every coaccessible $q \in P$, \ref{constr_sync:move_states:2} holds as well without modification, 
\ref{constr_sync:move_states:3} holds by taking
the set of all coaccessible states instead of $P$, and in \ref{constr_sync:move_states:4} we have to consider
only those maximal $[q]_{\eqreach}$ containing at least one final state.

Lastly, we state that for trim PDFAs, considering multiple starting
states does not change the complexity.

\begin{proposition} \label{constr_sync:prop:start_states}
	Let $\automat{B}{P}{\Sigma}{\mu}{p_0}{F}$ be a trim PDFA.	
	Then 
	\[ 
	\constrsync{L(\mathcal B)} \eqACZEROm \constrsync{L(\mathcal B) \cup \bigcup_{p \in I} L(\mathcal B_{p, F})}
	\]
	for all $I \subseteq P$.
\end{proposition}
\begin{proof}
	Let $u \in L(\mathcal B_{p, F})$ for some $p \in I$.
	As $\mathcal B$ is trim, there exists $v \in \Sigma^*$
	such that $\mu(p_0, v) = p$ and we get $vu \in L(\mathcal B)$.
	So $\bigcup_{p \in I} L(\mathcal B_{p, F}) \subseteq \suff{L(\mathcal B)}$
	and by \cref{constrsync:genresults:cor:subset} 
	the claim follows.
\end{proof}

\subsection{Concatenation and Complexity Lower Bounds}
\label{constrsync:sec:constraut:concatlowerbounds}

The following statement is a reformulation 
of \crefitem{constrsync:cor:UanL}{constrsync:cor:UanL:1} and 
\cref{constr_sync:prop:Gamma_suffix}
in terms of constraint PDFAs to obtain
complexity lower bounds.

\begin{proposition}
	Let $\mathcal B = (\Sigma, P, \mu, p_0, F)$
	be a PDFA and set
	\begin{enumerate}
		\item $\Gamma_1 = \{ x \in \Sigma \mid \applytrans{\mu}{p_0}{x}\text{ is undefined } \}$ and
		
		\item $\Gamma_2 = \{ x \in \Sigma \mid \forall p \in P \colon \applytrans{\mu}{p}{x} \notin F \}$.\footnote{We consider $\applytrans{\mu}{p}{x} \notin F$
			to be true also when $\applytrans{\mu}{p}{x}$ is undefined;
see also the convention from \cref{constr_sync:subsec:pfun} 
where for partial functions $f$, $f(x) \ne y$ included the case that $f(x)$ is undefined.}
	\end{enumerate}	
	If $U_1 \subseteq \Gamma_1^*$,
	$U_2 \subseteq \Gamma_2^*$ are regular,
	then $\constrsync{L(\mathcal B)} \leACZEROm \constrsync{U_1 L(\mathcal B) U_2}$.		
\end{proposition}
\begin{proof}
	As $\Gamma_1 \Sigma^* \cap L(\mathcal B) = \emptyset$
	and $\Sigma^* \Gamma_2 \cap L(\mathcal B) = \emptyset$
	the claim follows from  \crefitem{constrsync:cor:UanL}{constrsync:cor:UanL:1}
	and  \cref{constr_sync:prop:Gamma_suffix}.	
\end{proof}


\subsection{Equivalence to Unconstrained Synchronization} 
\label{constrsync:sec:constraut:unconstrained}

First, \cref{constrsync:prop:upperset} states
a rather general condition implying that the factor set
of the language accepted by a PDFA equals $\Sigma^*$.
Then, \cref{constrsync:cor:upperset}
derives several conditions implying equivalence
to $\constrsync{\Sigma^*}$.

\begin{proposition}\label{constrsync:prop:upperset}
	Let $\automat{B}{P}{\Sigma}{\mu}{p_0}{F}$ be a PDFA.
	Suppose there exists a set $S \subseteq P$ of coaccessible states
	containing an accessible state $s \in S$
	and	such that $\mu(p, x) \in S$ for each $p \in S$ and $x \in \Sigma$.
    Then $\fact{L(\mathcal B)} = \Sigma^*$.
\end{proposition} 
\begin{proof}
	Let $v \in \Sigma^*$. 
    As $s$ is accessible there exists $u \in \Sigma^*$
	such that $\mu(p_0, u) = s$.
	By assumption $\mu(p_0, uv) \in S$
	and there exists $w \in \Sigma^*$ such that $\mu(p_0, uvw) \in F$.
	So, $v \in \fact{L(\mathcal B)}$
	and we get $\fact{L(\mathcal B)} = \Sigma^*$.	
\end{proof} 

\begin{corollary}\label{constrsync:cor:upperset}
	Let $\automat{B}{P}{\Sigma}{\mu}{p_0}{F}$ be a PDFA.
	Then 
	$\constrsync{L(\mathcal B)} \eqACZEROm \constrsync{\Sigma^*}$
	and so $\constrsync{L(\mathcal B)}$ is \NL-complete
	for $|\Sigma| \ge 2$
	and \LOGSPACE-complete for $|\Sigma| = 1$
	in the following cases:
	\begin{enumerate}
		\item \label{constrsync:cor:upperset:item:S} $\mathcal B$ is trim and there exists a non-empty $S \subseteq P$ such that $\mu(p,x) \in S$ for all $p \in S$ and $x \in \Sigma$. 
		
		\item \label{constrsync:cor:upperset:item:complete} $\mathcal B$ is complete
		and every state is coaccessible.
		
		\item \label{constrsync:cor:upperset:item:trap} $\mathcal B$ contains an accessible trap state
		that is a final state, i.e., 
		a state $p \in F$ with $L(\mathcal B_{p,\{p\}}) = \Sigma^*$
		and $\mu(p_0, u) = p$ for some $u \in \Sigma^*$.
		
		\item \label{constrsync:cor:upperset:item:trimtrap} $\mathcal B$ is trim and has a trap state.     
	\end{enumerate}
\end{corollary} 
\begin{proof}
	In the following, we only point out how to obtain
	a set $S$ such that we can deduce $\fact{L(\mathcal B)} = \Sigma^*$
	from \cref{constrsync:prop:upperset}.
	Then, the claim follows
	from \cref{gen_results:cor:factors_etc_complexity}, 
    \cref{prelim:complexity:general_aut_sync,thm:SYNC-solv}.

\ref{constrsync:cor:upperset:item:S}: As $\mathcal B$ is trim, $S$ from the statement fulfills the conditions in \cref{constrsync:prop:upperset}.

\ref{constrsync:cor:upperset:item:complete}: Consider $S = P$ for \cref{constrsync:prop:upperset}.

\ref{constrsync:cor:upperset:item:trap}: Consider the set $S$ containing the accessible final trap state.
	
\ref{constrsync:cor:upperset:item:trimtrap}: A coaccessible trap state must be a final state.
 Thus, since~$\mathcal B$ is trim, each of its trap states is final
 and the claim follows from \cref{constrsync:cor:upperset:item:trap}.
\end{proof} 

\begin{corollary}
	Let $\automat{B}{P}{\Sigma}{\mu}{p_0}{F}$ be a trim PDFA.
    If for each $u\in \Sigma^*$ we have $\mu(p, u) = q$ for some $p,q \in P$,
    then  $\constrsync{L(\mathcal B)} \eqACZEROm \constrsync{\Sigma^*}$.
\end{corollary}
\begin{proof}
	Let $u \in \Sigma^*$. Then  $\mu(p, u) = q$ for some $p,q \in P$.
	Since~$\mathcal B$ is trim,
	we have $\mu(p_0, v) = p$
	and $\mu(q, w) \in F$ for some $v,w \in \Sigma^*$.
	Thus $u \in \factor{L(\mathcal B)}$,
	which gives $\factor{L(\mathcal B)} = \Sigma^*$.
	So, the claim follows
	with \cref{gen_results:cor:factors_etc_complexity}, 
	\cref{prelim:complexity:general_aut_sync,thm:SYNC-solv}.	
\end{proof}

\begin{remark}\label{rem:PDFA-for-constraint}
	For a trim automaton $\mathcal B$ with a totally defined transition function,
	by~\crefitem{constrsync:cor:upperset}{constrsync:cor:upperset:item:complete}
	we have $\constrsync{L(\mathcal B)} \eqACZEROm \constrsync{\Sigma^*}$.
	Hence, to describe constraint languages giving other complexities than the unconstrained problem,
	automata with totally defined transitions functions have useless states.
	So, out of this reason, we opted to use PDFAs to specify constraint languages
	as a more compact description than automata 
	with a totally defined transition function.
\end{remark}